%% Generated by tools/gen_error_codes.py from docs/errors/ of the compiler;
%% edit the compiler's pages or tools/error_themes.txt, then rerun it.

\clearpage
\section{Tuples, arrays, maps and options}
\label{sec:codes:compound}

The book's chapter \textit{Native compound types} specifies the rules behind these codes.

\begin{longtable}{@{}l p{0.78\linewidth} r@{}}
\toprule Code & Message & Page\\ \midrule \endhead
\hyperref[sec:codes:e4036]{E4036} & using a deep copy instead of a one level copy is prohibited when constructing a map \errhole{} & \pageref{sec:codes:e4036}\\
\hyperref[sec:codes:e4037]{E4037} & creating an instance of map \errhole{} on the stack is prohibited & \pageref{sec:codes:e4037}\\
\hyperref[sec:codes:e4121]{E4121} & the type \errhole{} cannot be used as a key, it is not comparable & \pageref{sec:codes:e4121}\\
\hyperref[sec:codes:e4122]{E4122} & the type \errhole{} cannot be used as a key, it is not hashable & \pageref{sec:codes:e4122}\\
\hyperref[sec:codes:e4131]{E4131} & mismatch tuple arity \errhole{} and \errhole{} & \pageref{sec:codes:e4131}\\
\hyperref[sec:codes:e4163]{E4163} & \errhole{} is not a tuple type & \pageref{sec:codes:e4163}\\
\hyperref[sec:codes:e4175]{E4175} & option of type \errhole{} has no error & \pageref{sec:codes:e4175}\\
\hyperref[sec:codes:e4176]{E4176} & option of type \errhole{} has no value & \pageref{sec:codes:e4176}\\
\hyperref[sec:codes:e4179]{E4179} & static array size \errhole{} exceeds limits \errhole{} & \pageref{sec:codes:e4179}\\
\hyperref[sec:codes:e4213]{E4213} & tuple left affectation is only defined for the operator \errhole{} & \pageref{sec:codes:e4213}\\
\hyperref[sec:codes:e4214]{E4214} & tuple access out of bound(\errhole{}), tuple arity is \errhole{} & \pageref{sec:codes:e4214}\\
\hyperref[sec:codes:e4267]{E4267} & operator \errhole{} is useless, the left operand of type \errhole{} always has a value & \pageref{sec:codes:e4267}\\
\bottomrule
\end{longtable}

\errorcode{E4036}{using a deep copy instead of a one level copy is prohibited when constructing a map \errhole{}}
\label{sec:codes:e4036}

Raised during \textbf{validation}, from \texttt{ValidateErrorMessage::\allowbreak{}CTOR\_\allowbreak{}MAP\_\allowbreak{}DCOPY} (\textit{src/\allowbreak{}ymirc/\allowbreak{}semantic/\allowbreak{}validator/\allowbreak{}errors.yr}).

\subsubsection*{What it means}

A map was constructed with \tokennolst{dcopy} where \tokennolst{copy} is required, cf. \hyperref[sec:codes:retired]{E4033}.

\subsubsection*{Example}

\textit{test\_\allowbreak{}resources/\allowbreak{}lit\_\allowbreak{}map/\allowbreak{}test12.yr}:

%% check: skip
\begin{lstlisting}[style=coloredverbatimError]
// a map is built one level deep, so a deep copy of a map literal is refused
fn main () {
    let _m_ = dcopy [1 => 2];
}
\end{lstlisting}

\begin{lstlisting}[style=bashVerb, breaklines=true, escapechar=~]
Error[E4036] : using a deep copy instead of a one level copy is prohibited when constructing a map mut [i32 => mut i32]
 --> test_resources/lit_map/test12.yr:(3,21)
 3  ~\lstbox{\textSFxi{}}~     let _m_ = dcopy [1 => 2];
    ~\lstbox{\textSFv{}}~                     ^
    ~\lstbox{\textSFxi{}}~
    = help: replace dcopy by the simple copy operator
    ~\lstbox{\textSFxi{}}~
 3  -     let _m_ = dcopy [1 => 2];
 3  +     let _m_ = copy [1 => 2];
    ~\lstbox{\textSFxi{}}~
    = when validating test12::main -- (test_resources/lit_map/test12.yr:(2,4))
    ~\lstbox{\textSFii{}}~~\lstbox{\textSFx{}}~~\lstbox{\textSFx{}}~~\lstbox{\textSFx{}}~~\lstbox{\textSFx{}}~~\lstbox{\textSFx{}}~~\lstbox{\textSFx{}}~~\lstbox{\textSFx{}}~
\end{lstlisting}

\subsubsection*{How to fix it}

Use \tokennolst{copy} to construct the map.

\errorcode{E4037}{creating an instance of map \errhole{} on the stack is prohibited}
\label{sec:codes:e4037}

Raised during \textbf{validation}, from \texttt{ValidateErrorMessage::\allowbreak{}CTOR\_\allowbreak{}MAP\_\allowbreak{}STACK} (\textit{src/\allowbreak{}ymirc/\allowbreak{}semantic/\allowbreak{}validator/\allowbreak{}errors.yr}).

\subsubsection*{What it means}

A map was created on the stack. Like classes, maps are reference types.

\subsubsection*{Example}

\textit{test\_\allowbreak{}resources/\allowbreak{}lit\_\allowbreak{}map/\allowbreak{}test3.yr}:

%% check: skip
\begin{lstlisting}[style=coloredverbatimError]
fn main () {
    let dmut _ = [1 => 2];
    let dmut _ = ["foo" => copy [1, 2, 3]];

    let dmut _ = ["foo" => [1, 2], "bar" => [1]];
    let dmut _ = ["foo" => 1, 2 => 3];
}
\end{lstlisting}

\begin{lstlisting}[style=bashVerb, breaklines=true, escapechar=~]
Error[E4037] : creating an instance of map mut [i32 => mut i32] on the stack is prohibited
 --> test_resources/lit_map/test3.yr:(2,18)
 2  ~\lstbox{\textSFxi{}}~     let dmut _ = [1 => 2];
    ~\lstbox{\textSFv{}}~                  ^
    ~\lstbox{\textSFxi{}}~
    = help: did you mean to enclose it within a copy?
    ~\lstbox{\textSFxi{}}~
 2  -     let dmut _ = [1 => 2];
 2  +     let dmut _ = copy [1 => 2];
    ~\lstbox{\textSFii{}}~~\lstbox{\textSFx{}}~~\lstbox{\textSFx{}}~~\lstbox{\textSFx{}}~~\lstbox{\textSFx{}}~~\lstbox{\textSFx{}}~~\lstbox{\textSFx{}}~~\lstbox{\textSFx{}}~
\end{lstlisting}

\subsubsection*{How to fix it}

Construct it with \tokennolst{copy}.

\errorcode{E4121}{the type \errhole{} cannot be used as a key, it is not comparable}
\label{sec:codes:e4121}

Raised during \textbf{validation}, from \texttt{ValidateErrorMessage::\allowbreak{}MAP\_\allowbreak{}KEY\_\allowbreak{}NOT\_\allowbreak{}COMPARABLE} (\textit{src/\allowbreak{}ymirc/\allowbreak{}semantic/\allowbreak{}validator/\allowbreak{}errors.yr}).

\subsubsection*{What it means}

A type is used as a map key and defines no comparison, so the map could not tell two keys apart.

\subsubsection*{Example}

\textit{test\_\allowbreak{}resources/\allowbreak{}diagnostics/\allowbreak{}test111.yr}:

%% check: skip
\begin{lstlisting}[style=coloredverbatimError]
// a class used as a map key, hashable but not comparable
class Key {
    pub self () {}
    impl Hashable;
}

fn main () {
    let _m_ : [&Key => i32] = copy [copy Key () => 1];
}
\end{lstlisting}

\begin{lstlisting}[style=bashVerb, breaklines=true, escapechar=~]
Error[E4121] : the type &(test111::Key) cannot be used as a key, it is not comparable
 --> test_resources/diagnostics/test111.yr:(2,7)
 2  ~\lstbox{\textSFxi{}}~ class Key {
    ~\lstbox{\textSFv{}}~       ^^^
 --> test_resources/diagnostics/test111.yr:(8,36)
 8  ~\lstbox{\textSFxi{}}~     let _m_ : [&Key => i32] = copy [copy Key () => 1];
    ~\lstbox{\textSFv{}}~                                    -  error: undefined operator == for types &(test111::Key) and &(test111::Key)
    ~\lstbox{\textSFxi{}}~                                    ~\lstbox{\textSFviii{}}~~\lstbox{\textSFx{}}~ error: class &(test111::Key) has no method named opEquals
    ~\lstbox{\textSFxi{}}~                                    ~\lstbox{\textSFii{}}~~\lstbox{\textSFx{}}~ error: class &(test111::Key) has no method named opCmp
    ~\lstbox{\textSFxi{}}~
    = when validating -- (test_resources/diagnostics/test111.yr:(8,36))
    = when validating test111::main -- (test_resources/diagnostics/test111.yr:(7,4))
    ~\lstbox{\textSFii{}}~~\lstbox{\textSFx{}}~~\lstbox{\textSFx{}}~~\lstbox{\textSFx{}}~~\lstbox{\textSFx{}}~~\lstbox{\textSFx{}}~~\lstbox{\textSFx{}}~~\lstbox{\textSFx{}}~
\end{lstlisting}

\subsubsection*{How to fix it}

Define the comparison operator on the type, or key the map on something that already has one.

\errorcode{E4122}{the type \errhole{} cannot be used as a key, it is not hashable}
\label{sec:codes:e4122}

Raised during \textbf{validation}, from \texttt{ValidateErrorMessage::\allowbreak{}MAP\_\allowbreak{}KEY\_\allowbreak{}NOT\_\allowbreak{}HASHABLE} (\textit{src/\allowbreak{}ymirc/\allowbreak{}semantic/\allowbreak{}validator/\allowbreak{}errors.yr}).

\subsubsection*{What it means}

A type is used as a map key and cannot be hashed.

\subsubsection*{Example}

\textit{test\_\allowbreak{}resources/\allowbreak{}diagnostics/\allowbreak{}test112.yr}:

%% check: skip
\begin{lstlisting}[style=coloredverbatimError]
// a class used as a map key without being hashable
class Key {
    pub self () {}
}

fn main () {
    let _m_ : [&Key => i32] = copy [copy Key () => 1];
}
\end{lstlisting}

\begin{lstlisting}[style=bashVerb, breaklines=true, escapechar=~]
Error[E4122] : the type &(test112::Key) cannot be used as a key, it is not hashable
 --> test_resources/diagnostics/test112.yr:(2,7)
 2  ~\lstbox{\textSFxi{}}~ class Key {
    ~\lstbox{\textSFv{}}~       ^^^  error: class type test112::Key does not implement trait core::types::Hashable
 --> test_resources/diagnostics/test112.yr:(7,36)
 7  ~\lstbox{\textSFxi{}}~     let _m_ : [&Key => i32] = copy [copy Key () => 1];
    ~\lstbox{\textSFv{}}~                                    -
    ~\lstbox{\textSFxi{}}~
    = when validating -- (test_resources/diagnostics/test112.yr:(7,36))
    = when validating test112::main -- (test_resources/diagnostics/test112.yr:(6,4))
    ~\lstbox{\textSFii{}}~~\lstbox{\textSFx{}}~~\lstbox{\textSFx{}}~~\lstbox{\textSFx{}}~~\lstbox{\textSFx{}}~~\lstbox{\textSFx{}}~~\lstbox{\textSFx{}}~~\lstbox{\textSFx{}}~
\end{lstlisting}

\subsubsection*{How to fix it}

Make the type hashable, or key the map on a value derived from it that already is.

\errorcode{E4131}{mismatch tuple arity \errhole{} and \errhole{}}
\label{sec:codes:e4131}

Raised during \textbf{validation}, from \texttt{ValidateErrorMessage::\allowbreak{}MISMATCH\_\allowbreak{}TUPLE\_\allowbreak{}ARITY} (\textit{src/\allowbreak{}ymirc/\allowbreak{}semantic/\allowbreak{}validator/\allowbreak{}errors.yr}).

\subsubsection*{What it means}

Two tuples meet where they have to agree and do not hold the same number of elements. Both arities are in the message.

\subsubsection*{Example}

\textit{test\_\allowbreak{}resources/\allowbreak{}lit\_\allowbreak{}tuples/\allowbreak{}test21.yr}:

%% check: skip
\begin{lstlisting}[style=coloredverbatimError]
fn foo ()-> (i32, i32);

fn main () {
    let mut a = 0, mut b = 0, mut c = 0;

    (a, b, c) = foo ();
}
\end{lstlisting}

\begin{lstlisting}[style=bashVerb, breaklines=true, escapechar=~]
Error[E4131] : mismatch tuple arity 3 and 2
 --> test_resources/lit_tuples/test21.yr:(6,5)
 6  ~\lstbox{\textSFxi{}}~     (a, b, c) = foo ();
    ~\lstbox{\textSFv{}}~     ^               ^
    ~\lstbox{\textSFii{}}~~\lstbox{\textSFx{}}~~\lstbox{\textSFx{}}~~\lstbox{\textSFx{}}~~\lstbox{\textSFx{}}~~\lstbox{\textSFx{}}~~\lstbox{\textSFx{}}~~\lstbox{\textSFx{}}~
\end{lstlisting}

\subsubsection*{How to fix it}

Make the two tuples the same length. A destructuring that binds fewer names than the tuple has also produces this.

\errorcode{E4163}{\errhole{} is not a tuple type}
\label{sec:codes:e4163}

Raised during \textbf{validation}, from \texttt{ValidateErrorMessage::\allowbreak{}NOT\_\allowbreak{}A\_\allowbreak{}TUPLE} (\textit{src/\allowbreak{}ymirc/\allowbreak{}semantic/\allowbreak{}validator/\allowbreak{}errors.yr}).

\subsubsection*{What it means}

An operation defined only on tuples was applied to something else.

\subsubsection*{Example}

\textit{test\_\allowbreak{}resources/\allowbreak{}lit\_\allowbreak{}tuples/\allowbreak{}test22.yr}:

%% check: skip
\begin{lstlisting}[style=coloredverbatimError]
fn main () {
    let mut a = 0, mut b = 0;

    (a, b) = 12;
}
\end{lstlisting}

\begin{lstlisting}[style=bashVerb, breaklines=true, escapechar=~]
Error[E4163] : i32 is not a tuple type
 --> test_resources/lit_tuples/test22.yr:(4,12)
 4  ~\lstbox{\textSFxi{}}~     (a, b) = 12;
    ~\lstbox{\textSFv{}}~            ^ ^^
    ~\lstbox{\textSFii{}}~~\lstbox{\textSFx{}}~~\lstbox{\textSFx{}}~~\lstbox{\textSFx{}}~~\lstbox{\textSFx{}}~~\lstbox{\textSFx{}}~~\lstbox{\textSFx{}}~~\lstbox{\textSFx{}}~
\end{lstlisting}

\subsubsection*{How to fix it}

Use a tuple value.

\errorcode{E4175}{option of type \errhole{} has no error}
\label{sec:codes:e4175}

Raised during \textbf{validation}, from \texttt{ValidateErrorMessage::\allowbreak{}OPTION\_\allowbreak{}HAS\_\allowbreak{}NO\_\allowbreak{}ERROR} (\textit{src/\allowbreak{}ymirc/\allowbreak{}semantic/\allowbreak{}validator/\allowbreak{}errors.yr}).

\subsubsection*{What it means}

The error side of an option was read on an option that never holds one.

\subsubsection*{Example}

\textit{test\_\allowbreak{}resources/\allowbreak{}diagnostics/\allowbreak{}test116.yr}:

%% check: skip
\begin{lstlisting}[style=coloredverbatimError]
// the error of an option known to hold a value
fn main () {
    let x : (i32)? = (1)?;
    let _ = x.error;
}
\end{lstlisting}

\begin{lstlisting}[style=bashVerb, breaklines=true, escapechar=~]
Error[E4175] : option of type (i32)? has no error
 --> test_resources/diagnostics/test116.yr:(4,14)
 4  ~\lstbox{\textSFxi{}}~     let _ = x.error;
    ~\lstbox{\textSFv{}}~              ^
    ~\lstbox{\textSFxi{}}~
    = when validating test116::main -- (test_resources/diagnostics/test116.yr:(2,4))
    ~\lstbox{\textSFii{}}~~\lstbox{\textSFx{}}~~\lstbox{\textSFx{}}~~\lstbox{\textSFx{}}~~\lstbox{\textSFx{}}~~\lstbox{\textSFx{}}~~\lstbox{\textSFx{}}~~\lstbox{\textSFx{}}~
\end{lstlisting}

\subsubsection*{How to fix it}

Read the value instead. An option built from a non-throwing computation has no error, cf. \hyperref[sec:codes:e4078]{E4078}.

\errorcode{E4176}{option of type \errhole{} has no value}
\label{sec:codes:e4176}

Raised during \textbf{validation}, from \texttt{ValidateErrorMessage::\allowbreak{}OPTION\_\allowbreak{}HAS\_\allowbreak{}NO\_\allowbreak{}VALUE} (\textit{src/\allowbreak{}ymirc/\allowbreak{}semantic/\allowbreak{}validator/\allowbreak{}errors.yr}).

\subsubsection*{What it means}

The value side of an option was read on an option that never holds one.

\subsubsection*{Example}

\textit{test\_\allowbreak{}resources/\allowbreak{}diagnostics/\allowbreak{}test117.yr}:

%% check: skip
\begin{lstlisting}[style=coloredverbatimError]
// the value of an option known to be empty
fn main () {
    let x : (i32)? = none;
    let _ = x.value;
}
\end{lstlisting}

\begin{lstlisting}[style=bashVerb, breaklines=true, escapechar=~]
Error[E4176] : option of type (i32)? has no value
 --> test_resources/diagnostics/test117.yr:(4,14)
 4  ~\lstbox{\textSFxi{}}~     let _ = x.value;
    ~\lstbox{\textSFv{}}~              ^
    ~\lstbox{\textSFxi{}}~
    = when validating test117::main -- (test_resources/diagnostics/test117.yr:(2,4))
    ~\lstbox{\textSFii{}}~~\lstbox{\textSFx{}}~~\lstbox{\textSFx{}}~~\lstbox{\textSFx{}}~~\lstbox{\textSFx{}}~~\lstbox{\textSFx{}}~~\lstbox{\textSFx{}}~~\lstbox{\textSFx{}}~
\end{lstlisting}

\subsubsection*{How to fix it}

Read the error instead, or build the option from a computation that can succeed.

\errorcode{E4179}{static array size \errhole{} exceeds limits \errhole{}}
\label{sec:codes:e4179}

Raised during \textbf{validation}, from \texttt{ValidateErrorMessage::\allowbreak{}OVERFLOW\_\allowbreak{}CAPACITY\_\allowbreak{}ARRAY} (\textit{src/\allowbreak{}ymirc/\allowbreak{}semantic/\allowbreak{}validator/\allowbreak{}errors.yr}).

\subsubsection*{What it means}

A static array is declared with a size beyond what the compiler accepts.

\subsubsection*{Example}

\textit{test\_\allowbreak{}resources/\allowbreak{}diagnostics/\allowbreak{}test118.yr}:

%% check: skip
\begin{lstlisting}[style=coloredverbatimError]
// a static array too large for the stack
fn main () {
    let _x_ : [i32 ; 100000000] = [0 ; 100000000];
}
\end{lstlisting}

\begin{lstlisting}[style=bashVerb, breaklines=true, escapechar=~]
Error[E4179] : static array size 100000000 exceeds limits 1048576
 --> test_resources/diagnostics/test118.yr:(3,40)
 3  ~\lstbox{\textSFxi{}}~     let _x_ : [i32 ; 100000000] = [0 ; 100000000];
    ~\lstbox{\textSFv{}}~                                        ^^^^^^^^^
    ~\lstbox{\textSFxi{}}~
    = when validating test118::main -- (test_resources/diagnostics/test118.yr:(2,4))
    ~\lstbox{\textSFii{}}~~\lstbox{\textSFx{}}~~\lstbox{\textSFx{}}~~\lstbox{\textSFx{}}~~\lstbox{\textSFx{}}~~\lstbox{\textSFx{}}~~\lstbox{\textSFx{}}~~\lstbox{\textSFx{}}~
\end{lstlisting}

\subsubsection*{How to fix it}

Reduce the size, or allocate a slice at run time instead.

\errorcode{E4213}{tuple left affectation is only defined for the operator \errhole{}}
\label{sec:codes:e4213}

Raised during \textbf{validation}, from \texttt{ValidateErrorMessage::\allowbreak{}TUPLE\_\allowbreak{}AFFECT\_\allowbreak{}OPERATOR} (\textit{src/\allowbreak{}ymirc/\allowbreak{}semantic/\allowbreak{}validator/\allowbreak{}errors.yr}).

\subsubsection*{What it means}

A tuple on the left of an assignment was used with an operator other than the plain one. Destructuring assignment is defined for \tokennolst{=} alone.

\subsubsection*{Example}

\textit{test\_\allowbreak{}resources/\allowbreak{}lit\_\allowbreak{}tuples/\allowbreak{}test23.yr}:

%% check: skip
\begin{lstlisting}[style=coloredverbatimError]
fn foo ()-> (i32, i32);

fn main () {
    let mut a = 0, mut b = 0;

    (a, b) += foo ();
}
\end{lstlisting}

\begin{lstlisting}[style=bashVerb, breaklines=true, escapechar=~]
Error[E4213] : tuple left affectation is only defined for the operator =
 --> test_resources/lit_tuples/test23.yr:(6,5)
 6  ~\lstbox{\textSFxi{}}~     (a, b) += foo ();
    ~\lstbox{\textSFv{}}~     ^      ^^
    ~\lstbox{\textSFii{}}~~\lstbox{\textSFx{}}~~\lstbox{\textSFx{}}~~\lstbox{\textSFx{}}~~\lstbox{\textSFx{}}~~\lstbox{\textSFx{}}~~\lstbox{\textSFx{}}~~\lstbox{\textSFx{}}~
\end{lstlisting}

\subsubsection*{How to fix it}

Split the compound assignment into a read and a write.

\errorcode{E4214}{tuple access out of bound(\errhole{}), tuple arity is \errhole{}}
\label{sec:codes:e4214}

Raised during \textbf{validation}, from \texttt{ValidateErrorMessage::\allowbreak{}TUPLE\_\allowbreak{}ARITY\_\allowbreak{}OVERFLOW} (\textit{src/\allowbreak{}ymirc/\allowbreak{}semantic/\allowbreak{}validator/\allowbreak{}errors.yr}).

\subsubsection*{What it means}

A tuple element was accessed beyond the arity of the tuple. Both the index and the arity are in the message.

\subsubsection*{Example}

\textit{test\_\allowbreak{}resources/\allowbreak{}diagnostics/\allowbreak{}test71.yr}:

%% check: skip
\begin{lstlisting}[style=coloredverbatimError]
fn main () {
    let t = (1, 2);
    let _ = t._5;
}
\end{lstlisting}

\begin{lstlisting}[style=bashVerb, breaklines=true, escapechar=~]
Error[E4214] : tuple access out of bound(5), tuple arity is 2
 --> test_resources/diagnostics/test71.yr:(3,14)
 3  ~\lstbox{\textSFxi{}}~     let _ = t._5;
    ~\lstbox{\textSFv{}}~              ^
    ~\lstbox{\textSFii{}}~~\lstbox{\textSFx{}}~~\lstbox{\textSFx{}}~~\lstbox{\textSFx{}}~~\lstbox{\textSFx{}}~~\lstbox{\textSFx{}}~~\lstbox{\textSFx{}}~~\lstbox{\textSFx{}}~
\end{lstlisting}

\subsubsection*{How to fix it}

Correct the index; tuple elements are numbered from zero.

\errorcode{E4267}{operator \errhole{} is useless, the left operand of type \errhole{} always has a value}
\label{sec:codes:e4267}

Raised during \textbf{validation}, from \texttt{ValidateErrorMessage::\allowbreak{}USELESS\_\allowbreak{}OPTION\_\allowbreak{}OR\_\allowbreak{}ELSE} (\textit{src/\allowbreak{}ymirc/\allowbreak{}semantic/\allowbreak{}validator/\allowbreak{}errors.yr}).

\subsubsection*{What it means}

An option fallback operator was applied to a value that always holds one, so the fallback can never run.

\subsubsection*{Example}

\textit{test\_\allowbreak{}resources/\allowbreak{}lit\_\allowbreak{}options/\allowbreak{}test21.yr}:

%% check: skip
\begin{lstlisting}[style=coloredverbatimError]
// the left operand always has a value, the right one would never be evaluated
fn main () {
    let _a_ : i32 = 12 ?? 42;
}
\end{lstlisting}

\begin{lstlisting}[style=bashVerb, breaklines=true, escapechar=~]
Error[E4267] : operator ?? is useless, the left operand of type i32 always has a value
 --> test_resources/lit_options/test21.yr:(3,24)
 3  ~\lstbox{\textSFxi{}}~     let _a_ : i32 = 12 ?? 42;
    ~\lstbox{\textSFv{}}~                        ^^
    ~\lstbox{\textSFii{}}~~\lstbox{\textSFx{}}~~\lstbox{\textSFx{}}~~\lstbox{\textSFx{}}~~\lstbox{\textSFx{}}~~\lstbox{\textSFx{}}~~\lstbox{\textSFx{}}~~\lstbox{\textSFx{}}~
\end{lstlisting}

\subsubsection*{How to fix it}

Remove the operator and use the value directly.
